\chapter{Open Systems}
\label{open}

\section{Variable mass and the Meshchersky equation}

Ordinary Newtonian mechanics implicitly assumes a fixed set of
particles. Rockets, raindrops growing by accretion, chains lifted
off a pile, and avalanches all violate that assumption: the mass $m$
of the body under study is itself a function of time. The correct
generalization of Newton's second law to such an \emph{open system}
--- worked out around the turn of the twentieth century by
I.\,V.\ Meshchersky --- is
\begin{equation}
  m\,\frac{d\mathbf v}{dt} \;=\; \mathbf F_{\mathrm{ext}}
  + \mathbf u\,\frac{dm}{dt} ,
  \label{meshchersky}
\end{equation}
where $\mathbf u$ is the velocity of the incoming or outgoing mass
\emph{relative to the body itself}, not relative to the lab frame.
Setting $\mathbf F_{\mathrm{ext}}=0$ and integrating
\eqref{meshchersky} with $\mathbf u$ constant and antiparallel to
the direction of travel yields the Tsiolkovsky rocket equation,
\begin{equation}
  v(t) - v(0) \;=\; u\,\ln\frac{m(0)}{m(t)} .
  \label{tsiolkovsky}
\end{equation}
The resemblance to Section~\ref{cm:hertz} is not superficial. The
term $\mathbf u\,\dot m$ in \eqref{meshchersky} looks, from the
point of view of the rocket body alone, exactly like an applied
force --- "thrust" --- and is routinely called one. But viewed as
part of the larger \emph{closed} system consisting of the rocket
\emph{together with} every parcel of exhaust it has ever ejected,
whose total momentum is exactly conserved, there is no force there
at all: only the ordinary transfer of momentum between two pieces of
one conserved whole. What looks like a force is, once again, an
artifact of insisting on an equation of motion for part of a system
rather than for the whole of it.

This is also a genuine pedagogical trap, worth flagging explicitly.
It is tempting to write $\mathbf F = d(m\mathbf v)/dt = m\dot{\mathbf
v} + \mathbf v\,\dot m$ and treat this directly as the equation of
motion of the open body. That expression is correct as a statement
about the rate of change of the momentum of the body \emph{alone},
but it is not, in general, equivalent to \eqref{meshchersky} and
gives the wrong dynamics unless the accreted or ejected mass happens
to leave (or join) with zero velocity relative to the body, i.e.\
$\mathbf u \equiv \mathbf v$ --- a special case, not the generic one,
and not the one realized by a rocket engine, whose whole purpose is
to eject mass at high relative velocity.

\section{Dissipative systems: putting friction back into a
  variational principle}

A second, quite different kind of "openness" arises whenever a
system exchanges energy with an environment that is not explicitly
modelled --- friction, drag, radiation reaction. None of these
forces derive from a potential in the ordinary sense, so the
Lagrangian and Hamiltonian formalisms, built for closed and
energy-conserving systems, do not apply directly. Two classical
devices repair this.

The first is Rayleigh's dissipation function $\mathcal
F_{\mathrm{diss}}(\dot q)$, quadratic in the generalized velocities
for linear damping, which modifies the Euler--Lagrange equations to
\begin{equation}
  \frac{d}{dt}\frac{\partial L}{\partial \dot q^i}
  - \frac{\partial L}{\partial q^i}
  \;=\; -\,\frac{\partial \mathcal F_{\mathrm{diss}}}{\partial \dot
  q^i}.
  \label{rayleigh}
\end{equation}
This correctly reproduces the (irreversible) equations of motion,
but it is a hybrid: $\mathcal F_{\mathrm{diss}}$ is not itself part
of a genuine, self-contained stationary-action principle for $L$
alone.

The second, due to Bateman (1931), is more radical and more
illuminating: a fully conservative, symplectic Lagrangian \emph{can}
be written for a damped oscillator, but only at the price of
doubling the number of degrees of freedom. Introduce an auxiliary
"mirror" coordinate $y$ alongside the physical coordinate $x$,
\begin{equation}
  L_{\mathrm{dual}}(x,\dot x,y,\dot y) \;=\;
  m\,\dot x\,\dot y \;+\; \frac{\gamma}{2}\bigl(x\dot y - \dot x
  y\bigr) \;-\; k\,xy ,
  \label{bateman}
\end{equation}
whose Euler--Lagrange equations are precisely the damped oscillator,
$\ddot x + (\gamma/m)\dot x + \omega_0^2 x = 0$, for the physical
coordinate $x$, paired with a time-reversed, exponentially
\emph{amplified} equation for $y$. The energy the
physical oscillator loses is exactly the energy its mirror partner
gains, so the doubled system is conservative even though either half
alone is not. This is a small, entirely classical and elementary
demonstration of an idea that later became central to the theory of
open quantum systems: dissipation and irreversibility in a subsystem
are not exceptions to an underlying conservative dynamics, but the
generic appearance of a conservative dynamics once part of the full
system --- here, quite literally, the mirror half --- has been
discarded from the description.

\section{When "open" turns pathological: radiation reaction}

Not every attempt to write a closed-looking equation for an open
system ends benignly. An accelerating point charge radiates, and by
momentum conservation it must experience a recoil, or self-force, in
consequence. Treating the charge \emph{in isolation} --- an open
system in the strict sense, since the charge together with its own
radiated field is really a system with infinitely many degrees of
freedom, formally "integrated out" down to a single effective force
term --- leads to the (nonrelativistic) Abraham--Lorentz equation,
\begin{equation}
  m\,\ddot{\mathbf x} \;=\; \mathbf F_{\mathrm{ext}} \;+\;
  m\tau\,\dddot{\mathbf x} ,
  \qquad
  \tau \;=\; \frac{q^2}{6\pi\varepsilon_0\,m\,c^3}
  \;=\; \frac{2}{3}\,\frac{r_e}{c},
  \label{abraham-lorentz}
\end{equation}
where $r_e$ is the classical electron radius. Because
\eqref{abraham-lorentz} is third order in position, it requires
more initial data than position and velocity to determine a unique
solution, and generic initial data produce \emph{runaway} solutions
in which the self-acceleration grows without bound even with
$\mathbf F_{\mathrm{ext}}=0$. Suppressing the runaway by hand forces
the opposite pathology: the charge must begin accelerating
\emph{before} an external force is switched on --- a small, short-lived,
but genuinely acausal-looking pre-acceleration. (In practice, both
pathologies signal that \eqref{abraham-lorentz} is being pushed
outside its domain of validity, on timescales comparable to $\tau
\sim 10^{-24}\,\mathrm s$; treated properly, as an effective equation
valid only for slowly varying external forces --- via, e.g., the
Landau--Lifshitz reduction of order --- the pathologies disappear.)
The lesson is a cautionary complement to the rest of this section:
"integrating out" the unmodelled part of an open system is not
always harmless bookkeeping. Sometimes it is, as with the Meshchersky
equation; sometimes, as here, it produces an equation that is simply
malformed unless handled with real care.

